
\documentclass[11pt]{article}
\usepackage[a4paper,margin=1in]{geometry}
\usepackage{graphicx}
\usepackage{booktabs}
\usepackage{hyperref}
\usepackage{longtable}
\usepackage{amsmath}
\usepackage{float}
\title{Production Analysis of Blockade-Scale-Normalised KAN-Jastrow and Graph-KAN Neural Quantum States for Rydberg AMO Benchmarks}
\author{GraphKAN AMO NQS production artifact}
\date{20 June 2026}
\begin{document}
\maketitle
\begin{abstract}
This report audits the repaired proposal and the uploaded prototype code, implements a hardware-free CPU production-pilot engine, and records the experiments that could be executed in the current environment. The artifact validates Rydberg Hamiltonian construction and exact diagonalisation on small systems, runs exact full-basis variational pilots for product, Jastrow, MLP-pair random-feature and KAN-Jastrow models, extracts gauge-fixed learned radial functions, and generates reproducible figures and tables. It does not claim top-journal readiness, first KAN-NQS, Graph-KAN value, quantum advantage, or hardware validation.
\end{abstract}

\section{Executive summary}
The result is a real CPU benchmark engine, not a proposal-only scaffold. The available environment lacked NetKet, Flax, QuSpin and TeNPy, so production VMC and DMRG profiles are recorded as not executed. The feasible path produced exact references and exact-fullsum pilot evidence. The correct classification is pilot benchmark continuation, not paper-ready PRA or Quantum evidence.

\section{Code and proposal audit}
The uploaded proposal was already anti-scoop repaired: it forbids first KAN-NQS and KAN-superiority claims and centres the project on blockade-scale normalisation, gauge fixing, and fair Rydberg benchmarks. The uploaded codebase passed its original tests and smoke script but was a PyTorch smoke scaffold. The present repository replaces the scientific core with deterministic geometry, Hamiltonian, ED, exact-fullsum, statistics and reporting modules.

\section{Software environment}
The release manifest records Python, NumPy, SciPy, pandas and optional dependency status. Optional NetKet/JAX/Flax VMC remains a required next step for a final paper. The CPU implementation is intentionally hardware-free.

\section{Hamiltonian validation}
The Rydberg Hamiltonian used here is
\[
H/\hbar=\sum_i \frac{\Omega_i}{2}\sigma_i^x-\sum_i\Delta_i n_i+\sum_{i<j}\frac{C_6}{r_{ij}^6}n_i n_j.
\]
Dense and sparse builders are checked against each other, the $\Omega=0$ diagonal limit is checked analytically, and exact references are generated where feasible.

\begin{figure}[H]\centering\includegraphics[width=.75\linewidth]{../figures/fig_ed_validation.pdf}\caption{Hamiltonian validation from generated reference files.}\end{figure}

\section{Exact-reference generation}
References were generated for the configured chain, square and ladder instances. Exact diagonalisation uses dense diagonalisation below a dimension threshold and SciPy sparse ARPACK above it. Large-system truth is not claimed.

\section{Model and baseline definitions}
The CPU pilot uses linear-in-parameters exact-basis feature models: product, radial/smooth Jastrow, MLP-pair random features with the same blockade-normalised inputs as KAN-Jastrow, and KAN-Jastrow radial basis features. This is sufficient for Hamiltonian/model plumbing and pilot comparisons, but it is not a substitute for final NetKet/Flax trainable baselines.

\section{Experimental profiles executed}
Executed profiles are quick CI, ED references, exact full-sum pilot, Rydberg chain pilot, aggregation, figures, tables, report, manifest, claim validation and readiness checks. PRA-core, geometry transfer, Graph-KAN, robustness, Bose-Hubbard and IBM profiles are not executed in this CPU artifact.

\section{Results}
\subsection{Energy and correlations}
\begin{center}
\resizebox{\linewidth}{!}{%
\begin{tabular}{lllrrrr}
\toprule
profile & model & metric & median & ci\_low & ci\_high & n \\
\midrule
exact\_fullsum\_pilot & kan\_jastrow & correlation\_mae & 0.014202 & 0.009446 & 0.015588 & 54 \\
exact\_fullsum\_pilot & kan\_jastrow & density\_mae & 0.042149 & 0.028581 & 0.055463 & 54 \\
exact\_fullsum\_pilot & kan\_jastrow & energy\_density\_error & 0.016203 & 0.013918 & 0.019963 & 54 \\
exact\_fullsum\_pilot & kan\_jastrow & fidelity & 0.944592 & 0.940048 & 0.955879 & 54 \\
exact\_fullsum\_pilot & kan\_jastrow & parameter\_count & 13.000000 & 13.000000 & 13.000000 & 54 \\
exact\_fullsum\_pilot & kan\_jastrow & runtime\_sec & 0.002353 & 0.001835 & 0.002905 & 54 \\
exact\_fullsum\_pilot & mlp\_pair & correlation\_mae & 0.015460 & 0.009440 & 0.019192 & 54 \\
exact\_fullsum\_pilot & mlp\_pair & density\_mae & 0.057710 & 0.033198 & 0.083868 & 54 \\
exact\_fullsum\_pilot & mlp\_pair & energy\_density\_error & 0.016975 & 0.015284 & 0.019963 & 54 \\
exact\_fullsum\_pilot & mlp\_pair & fidelity & 0.934155 & 0.899831 & 0.953699 & 54 \\
exact\_fullsum\_pilot & mlp\_pair & parameter\_count & 12.000000 & 12.000000 & 12.000000 & 54 \\
exact\_fullsum\_pilot & mlp\_pair & runtime\_sec & 0.003952 & 0.003281 & 0.004720 & 54 \\
exact\_fullsum\_pilot & product & correlation\_mae & 0.025977 & 0.021337 & 0.032700 & 54 \\
exact\_fullsum\_pilot & product & density\_mae & 0.160077 & 0.112417 & 0.188048 & 54 \\
exact\_fullsum\_pilot & product & energy\_density\_error & 0.111226 & 0.071524 & 0.151655 & 54 \\
exact\_fullsum\_pilot & product & fidelity & 0.581838 & 0.429031 & 0.724609 & 54 \\
exact\_fullsum\_pilot & product & parameter\_count & 2.000000 & 2.000000 & 2.000000 & 54 \\
exact\_fullsum\_pilot & product & runtime\_sec & 0.000863 & 0.000783 & 0.000938 & 54 \\
exact\_fullsum\_pilot & radial\_jastrow & correlation\_mae & 0.014141 & 0.009493 & 0.015497 & 54 \\
exact\_fullsum\_pilot & radial\_jastrow & density\_mae & 0.037919 & 0.028581 & 0.053390 & 54 \\
\bottomrule
\end{tabular}

}
\end{center}
\begin{figure}[H]\centering\includegraphics[width=.8\linewidth]{../figures/fig_pilot_energy_correlation.pdf}\caption{Pilot median energy-density and connected-correlation errors.}\end{figure}

\subsection{Paired statistics}
\begin{center}
\resizebox{\linewidth}{!}{%
\begin{tabular}{llrrrrr}
\toprule
comparison & metric & median\_delta\_baseline\_minus\_candidate & ci\_low & ci\_high & n\_pairs & p\_value \\
\midrule
mlp\_pair\_minus\_kan\_jastrow & energy\_density\_error & 0.000000 & 0.000000 & 0.000026 & 71 & 0.000005 \\
product\_minus\_kan\_jastrow & energy\_density\_error & 0.058024 & 0.031446 & 0.111529 & 71 & 0.000000 \\
radial\_jastrow\_minus\_kan\_jastrow & energy\_density\_error & 0.000000 & 0.000000 & 0.000000 & 55 & 0.000476 \\
smooth\_jastrow\_minus\_kan\_jastrow & energy\_density\_error & -0.000275 & -0.000945 & 0.001020 & 16 & 1.000000 \\
mlp\_pair\_minus\_kan\_jastrow & correlation\_mae & -0.000000 & -0.000001 & 0.000000 & 71 & 0.578352 \\
product\_minus\_kan\_jastrow & correlation\_mae & 0.007069 & 0.000568 & 0.009494 & 71 & 0.000000 \\
radial\_jastrow\_minus\_kan\_jastrow & correlation\_mae & 0.000000 & -0.000000 & 0.000000 & 55 & 0.814520 \\
smooth\_jastrow\_minus\_kan\_jastrow & correlation\_mae & -0.000367 & -0.005373 & -0.000003 & 16 & 0.129730 \\
\bottomrule
\end{tabular}

}
\end{center}
\begin{figure}[H]\centering\includegraphics[width=.8\linewidth]{../figures/fig_kan_vs_mlp_pair_ablation.pdf}\caption{Paired pilot ablations. Positive deltas favour KAN-Jastrow over the named baseline.}\end{figure}

\subsection{Compute}
\begin{figure}[H]\centering\includegraphics[width=.75\linewidth]{../figures/fig_compute_pareto.pdf}\caption{Pilot compute-error Pareto plot.}\end{figure}

\section{Interpretability analysis}
Gauge fixing subtracts a constant offset, centres the learned pair function, normalises amplitude, and aligns a sign convention. This prevents raw redundant parameterisations from being interpreted as physical functions.
\begin{figure}[H]\centering\includegraphics[width=.8\linewidth]{../figures/fig_gauge_fixed_radial_functions.pdf}\caption{Gauge-fixed KAN-Jastrow radial functions from the CPU pilot.}\end{figure}
\begin{figure}[H]\centering\includegraphics[width=.8\linewidth]{../figures/fig_scale_collapse.pdf}\caption{Pilot scale-collapse diagnostic plotted against $r/R_\Omega$.}\end{figure}

\section{Ablations, robustness and failed runs}
The CPU pilot includes the central KAN-vs-MLP-pair ablation under matched blockade-normalised input features. Full basis-family, edge-cutoff, geometry-transfer and sampler robustness studies remain next experiments. Failed or unexecuted runs are retained rather than hidden.
\begin{center}
\resizebox{\linewidth}{!}{%
\begin{tabular}{lllrlll}
\toprule
profile & instance\_id & model & seed & stage & run\_status & reason \\
\midrule
rydberg\_pra\_core & profile & all & -1 & profile & failed & Requires NetKet/Flax and longer wall-clock budget; CPU artifact records command but does not fake results. \\
rydberg\_geometry\_transfer & profile & all & -1 & profile & failed & Requires trained parameter-conditioned models beyond CPU pilot. \\
graphkan\_full & profile & all & -1 & profile & failed & Graph-KAN is gated on KAN-Jastrow stability and a fair graph MPNN baseline. \\
ablation\_core & profile & all & -1 & profile & failed & Ablations are generated from pilot pair comparisons in this CPU artifact; full basis-family sweep is deferred. \\
robustness\_core & profile & all & -1 & profile & failed & Position jitter and sampler robustness require expanded runs beyond the CPU artifact. \\
bose\_hubbard\_optional & profile & all & -1 & profile & failed & Bose-Hubbard is secondary and should not delay Rydberg core. Tiny analytic tests are included in pytest. \\
\bottomrule
\end{tabular}

}
\end{center}

\section{Decision gates and venue readiness}
\begin{center}
\resizebox{\linewidth}{!}{%
\begin{tabular}{lll}
\toprule
gate & status & evidence \\
\midrule
A\_hamiltonian\_validity & PASS & Dense/sparse and analytic-limit tests passed in pytest and quick\_ci. \\
B\_kan\_jastrow\_viability & PILOT\_INCONCLUSIVE\_OR\_NEGATIVE & KAN-Jastrow is not convincingly better than matched MLP-pair/smooth Jastrow in the CPU pilot; method-superiority claims are unsafe. \\
C\_interpretability & PILOT\_PASS & Median pilot curve alignment = 0.999. \\
D\_graphkan\_value & NOT\_RUN & Graph-KAN profile is gated on KAN-Jastrow stability and was not executed. \\
E\_top\_journal\_pursuit & NOT\_READY & Current artifact is a CPU pilot, not PRA/Quantum-grade final evidence. \\
\bottomrule
\end{tabular}

}
\end{center}
\begin{figure}[H]\centering\includegraphics[width=.8\linewidth]{../figures/fig_claim_ladder_decision.pdf}\caption{Claim ladder and readiness gates.}\end{figure}

\section{What went well}
The Rydberg Hamiltonian path is testable and reproducible. The exact-fullsum engine makes model evaluation noise-free on small systems. The output schemas are long-format and figure/table generation is script-driven. Gauge-fixed learned functions are exported as data rather than manually drawn plots.

\section{What went badly}
The environment did not include NetKet, Flax, QuSpin or TeNPy. Therefore no production VMC, QuSpin cross-check, or DMRG reference was run. The CPU MLP and RBM baselines are pilot substitutes, not final strong baselines. Graph-KAN was not run because its value should be gated on KAN-Jastrow stability and a fair graph-MPNN comparison.

\section{What is inconclusive}
KAN-Jastrow viability remains inconclusive until more seeds, full trainable baselines, larger parameter sweeps, DMRG references for chains, and NetKet VMC are run. Interpretability is also pilot-level: a stable-looking curve is not enough without broad seed and parameter stability.

\section{Pursue or pivot}
The project is worth pursuing as a rigorous benchmark/software paper if the next phase confirms Hamiltonian validation, baseline fairness, and learned-function stability. It is not yet a method-superiority paper. If KAN-Jastrow is dominated by MLP-pair and smooth Jastrow after full fair comparison, pivot to a negative-results benchmark.

\section{Safe and unsafe claims}
Safe claims: the artifact implements and tests a hardware-free exact Rydberg pilot; the proposal's anti-scoop framing remains necessary; optional hardware is unnecessary for main claims. Unsafe claims: first KAN-NQS, general KAN superiority, Graph-KAN value, quantum advantage, IBM validation of Rydberg NQS, and top-journal readiness.

\section{Required next experiments}
Run NetKet/Flax VMC, QuSpin cross-checks, TeNPy DMRG for 1D, five to ten seeds, full RBM/MLP/MPNN baselines, geometry transfer, Graph-KAN only after KAN-Jastrow stability, and robustness/failure maps.

\begin{figure}[H]\centering\includegraphics[width=.8\linewidth]{../figures/fig_failure_map.pdf}\caption{Pilot failure/best-model map from actual result rows.}\end{figure}
\begin{figure}[H]\centering\includegraphics[width=.8\linewidth]{../figures/fig_reproducibility_workflow.pdf}\caption{Hardware-free reproducibility workflow.}\end{figure}

\section{Conclusion}
The artifact moves the project from Level 1 smoke scaffold to a Level 2 CPU exact-reference pilot engine. It is production-coding useful and scientifically honest, but not journal-ready. The next gate is a NetKet/Flax-backed Rydberg chain benchmark with strong trainable baselines and DMRG cross-checks.
\end{document}
